% answers_short.tex —— 2009 年·全国大纲卷Ⅱ·文科 · 简短答案速查
% 由 tools/make_short_answers.py 自 answers.tex 的 \gksolhead/\gkanswer 提取，
% 并与 2009/2009_数学真题_全国2文_答案要点.md 逐题比对后生成，勿手改（改源 MD 后重跑）。

\begin{gknotice}
\textbf{使用说明}：本节只给最终答案便于快速核对，完整推导见解析版。答案为编写者独立推导，\textbf{非官方标准答案}。
\end{gknotice}

\gkpart{选择题}
\begin{center}\begin{tabular}{ccccccccccccc}
\toprule
题号 & 1 & 2 & 3 & 4 & 5 & 6 & 7 & 8 & 9 & 10 & 11 & 12 \\
\midrule
答案 & C & D & A & D & C & C & 待核 & A & D & C & D & A \\
\bottomrule
\end{tabular}\end{center}

\gkpart{填空题}
\begin{itemize}[leftmargin=2.6em]
  \item[\textbf{13.}] $3$
  \item[\textbf{14.}] $6$
  \item[\textbf{15.}] $\dfrac{25}{4}$
  \item[\textbf{16.}] $8$
\end{itemize}

\gkpart{解答题}
\begin{itemize}[leftmargin=2.6em]
  \item[\textbf{17.}] $S_n=n(n-9)$ 或 $S_n=n(9-n)$。
  \item[\textbf{18.}] $B=60^\circ$。
  \item[\textbf{19.}] (1) 见解析版；(2) $B_1C$ 与平面 $BCD$ 所成的角为 $30^\circ$
  \item[\textbf{20.}] (1) 甲组、乙组各抽取 $2$ 人；(2) $\dfrac8{15}$；(3) $\dfrac{104}{225}$
  \item[\textbf{21.}] (1) 当 $x<2$ 或 $x>2a$ 时，$f(x)$ 单调递增；当 $2<x<2a$ 时，$f(x)$ 单调递减；(2) $1<a<6$
  \item[\textbf{22.}] (1) $a=\sqrt3,\ b=\sqrt2$；(2) 存在 $P\left(\dfrac32,\dfrac{\sqrt2}{2}\right)$、$P\left(\dfrac32,-\dfrac{\sqrt2}{2}\right)$；相应直线为 $y=\sqrt2(x-1)$、$y=-\sqrt2(x-1)$
\end{itemize}
